\cleardoubleoddpage
\chapter{Method}
\label{chap:method}

\section{Pipeline Overview}

\autoref{fig:pipeline} summarizes the proposed two-stage surrogate. The encoder is trained first and
then frozen. It converts every trajectory into a latent cache. Sequence models are trained only on
windows from this cache. During evaluation, predicted latents are passed through the frozen
diagnostic heads so that temporal errors can be measured in flux and spectra space.

\begin{figure}[htbp]
  \centering
  \includegraphics[width=0.95\textwidth]{pipeline_diagram.png}
  \caption{Latent-surrogate pipeline. A five-dimensional snapshot is embedded into a compact
  vector, evolved autoregressively, and evaluated through latent and diagnostic metrics.}
  \label{fig:pipeline}
\end{figure}

This staged design prevents file layout and spatial tensor size from leaking into the sequence
model. It also guarantees that all temporal baselines in a comparison use the same representation.

\section{Data Contract and Splitting}

A single snapshot has shape
\begin{equation}
  x_t \in \mathbb{R}^{C\times S_1\times S_2\times S_3\times S_4\times S_5},
\end{equation}
and a batch has shape
\begin{equation}
  X \in \mathbb{R}^{B\times C\times S_1\times S_2\times S_3\times S_4\times S_5}.
\end{equation}
The validated Cyclone snapshot shape is $(4,32,8,16,85,32)$ before batching. All adapters convert
samples to \code{float32} at the model boundary. The channel-first contract is validated rather
than inferred, because an unnoticed axis permutation would still be numerically executable but
scientifically invalid.

The Cyclone adapter loads a field snapshot at raw index $r$ and attaches flux and spectra from
$r+b$, where the configured bundle offset is $b=1$. Retained field snapshots are separated by the
subsampling stride of eight raw indices. Consequently, the diagnostic head learns
$z(r)\mapsto y(r+1)$, and the diagnostic reported for rollout horizon $h$ corresponds to the target
attached to the predicted retained state, $y(r+8h+1)$. Throughout the thesis, ``flux at horizon
$h$'' refers to this attached diagnostic target rather than a value at exactly the field index.

\begin{figure}[htbp]
  \centering
  \includegraphics[width=0.95\textwidth]{data_contract_diagram.png}
  \caption{Data and latent-cache contracts. Splits are defined at trajectory level so overlapping
  windows from one trajectory cannot appear in both training and evaluation.}
  \label{fig:data-contract}
\end{figure}

Trajectories are deterministically shuffled using a fixed seed and assigned with ratios
$0.8/0.1/0.1$. For the medium cache this produces 41 training, five validation, and five test
trajectories. Each trajectory contains 23 retained time steps, giving 1,173 latent vectors. Splitting
before window construction is essential: randomly splitting overlapping windows would leak nearly
identical temporal contexts across partitions.

For context length $T_c$ and prediction length $T_p$, a valid window begins at any index
$i\in\{0,\ldots,T-T_c-T_p\}$. The principal sequence experiment uses $T_c=8$ for one-step training
and recursive evaluation over $T_r=8$ forecast steps.

\section{Snapshot Encoder}

The selected encoder is a five-dimensional pointwise convolutional hierarchy. The input channel
axis is moved to the end for Flax operations. Each of three stages first subsamples the five spatial
axes and then applies a learned pointwise projection with 8, 16, and 32 output channels. The stride
vectors are
\begin{align}
  (2,2,2,4,2),\quad (2,1,2,4,2),\quad (2,1,1,2,2).
\end{align}
GELU activations follow each projection. Global mean pooling over all remaining spatial positions
produces one vector, and a dense layer maps it to $z_t\in\mathbb{R}^{128}$.

The $1^5$ kernels make this architecture computationally tractable for the large input tensor. Its
spatial mixing comes primarily from subsampling and global pooling rather than wide convolutional
kernels. This is an engineering baseline, not a claim that pointwise operations optimally represent
gyrokinetic geometry.

\section{Representation Objective}

Two augmented views are sampled independently from a snapshot. The enabled transforms are Gaussian
noise with relative scale $0.01$, amplitude jitter of $0.02$, and element masking with probability
$0.02$. Let $g$ be a two-layer projection head and $h$ a prediction head. The symmetric SimSiam loss
is
\begin{equation}
  \mathcal{L}_{\mathrm{Siam}}
  =\frac{1}{2}\left[-\cos\left(h(g(f(x^{(1)}))),\operatorname{sg}(g(f(x^{(2)})))\right)
  -\cos\left(h(g(f(x^{(2)}))),\operatorname{sg}(g(f(x^{(1)})))\right)\right],
\end{equation}
where $\operatorname{sg}$ denotes stop-gradient.

The latent vector is also passed to independent diagnostic heads. Each head contains a 256-unit
hidden layer and predicts scalar flux or a 32-bin spectrum. With mean squared diagnostic losses,
the encoder objective is
\begin{equation}
  \mathcal{L}_{\mathrm{enc}}
  = \mathcal{L}_{\mathrm{flux}}
  + \mathcal{L}_{\mathrm{spectra}}
  + 0.5\,\mathcal{L}_{\mathrm{Siam}}.
\end{equation}
Both spectra, \code{kyspec} and \code{fluxspec}, are used in the spectra term. The diagnostic heads
make the representation task-aware, while the Siamese term encourages consistency under the small
perturbations shown in \autoref{fig:encoder-objective}.

\begin{figure}[htbp]
  \centering
  \includegraphics[width=0.92\textwidth]{encoder_objective_diagram.png}
  \caption{Joint encoder training objective: two conservative augmented views are aligned with a
  stop-gradient Siamese objective, while flux and spectra heads retain diagnostic information.}
  \label{fig:encoder-objective}
\end{figure}

\section{Latent Cache}

After encoder training, every trajectory is embedded once. The HDF5 cache stores $z_t$, timestep
indices, physical time when available, flux, spectra, the encoder checkpoint path, and the resolved
configuration. This immutable boundary has two benefits: temporal models can be iterated without
re-reading five-dimensional fields, and all compared sequence models receive identical inputs.

For cache normalization, the mean $\mu_z$ and standard deviation $\sigma_z$ are computed from the
training trajectories only. Sequence learning then uses
\begin{equation}
  \tilde z_t = \frac{z_t-\mu_z}{\sigma_z+\epsilon},\qquad \epsilon=10^{-6}.
\end{equation}
Predictions are transformed back before the frozen diagnostic heads are evaluated. Restricting the
statistics to the training split prevents test information from entering preprocessing.

\section{Sequence Models}

\subsection{Persistence}

The persistence baseline predicts
\begin{equation}
  \hat z_{t+1}=z_t
\end{equation}
and repeats this value during recursive rollout. It is difficult to beat when adjacent states change
slowly, making it a necessary baseline for short-horizon dynamics.

\subsection{GRU}

The recurrent baseline uses one GRU layer with hidden width 256 and a dense output projection to
128 dimensions. It consumes the eight context latents sequentially and predicts the next latent
from the final hidden state.

\subsection{Causal Transformer}

The selected Transformer projects each 128-dimensional latent to width 256, adds learned positional
embeddings, and applies three pre-normalized causal self-attention blocks. Each block has four
attention heads and an MLP expansion ratio of four. The normalized final context token is projected
back to 128 dimensions. Dropout is zero in the medium experiment.

All learned sequence models minimize a latent objective combining mean squared error and cosine
dissimilarity. Adam-style optimization uses learning rate $10^{-3}$, weight decay $10^{-4}$, a
25-step warmup, and gradient-norm clipping at one. Models are trained for 500 update steps.

\section{Autoregressive Evaluation}

Given true context $z_{t-T_c+1:t}$, the model predicts $\hat z_{t+1}$. The context window is shifted,
the prediction is appended, and the procedure is repeated for eight steps. The persistence baseline
uses the same window construction. For latent prediction $\hat Z$ and target $Z$, the principal
latent metric is
\begin{equation}
  \mathrm{MSE}_{z}=\frac{1}{BT_rD}\lVert \hat Z-Z\rVert_2^2.
\end{equation}

Flux RMSE is computed after applying the frozen flux head $d_f$:
\begin{equation}
  \mathrm{RMSE}_{f}=\sqrt{\frac{1}{BT_r}\sum_{b=1}^{B}\sum_{k=1}^{T_r}
  \left(d_f(\hat z_{b,k})-f_{b,k}\right)^2}.
\end{equation}
Spectra MSE, log-MSE, relative $L_2$, and shape correlation are reported as secondary diagnostics.
A rollout is stable when every predicted value is finite.

Overlapping windows can cause long trajectories to dominate a sample-weighted aggregate. The
verified evaluation first computes metrics per trajectory and then gives each test trajectory equal
weight. Dispersion therefore represents between-trajectory variation. Horizons are one-based in
reports, so horizon one denotes the first predicted state.
